\section{A grading and a shift in genus zero}
\label{shift:section}

We now construct two commuting operators on the equivariant quantum
cohomology of $W$.  One is the grading operator.  The other translates
$\lambda$ by the loop variable $z$ and is defined by genus-zero invariants
of a bundle over $\mathbb P^1$.  Its fixed-locus expression will connect the
ordinary theories of $B$ and $X$.  Its spectrum, computed below using only
fiber curves, will provide the branches along which we transport the
constraints.

Throughout this section, $\star$ means the small equivariant quantum
product of $W$, at primary background zero.  Put
\[
 c=c_1^T(TW),\qquad
 \kappa=c_1(TB)+c_1(E).
\]
The equivariant projective-bundle formula gives
\begin{equation}
 c=(r+1)p+\pi^*\kappa+\lambda,
 \qquad
 \int_{(\beta,l)}c_1(TW)=(r+1)l+\int_\beta\kappa.
 \label{shift:c1}
\end{equation}
Choose a homogeneous Hodge basis of $H^*(B)$ and the resulting polynomial
basis $\pi^*b\,p^j$, $0\leq j\leq r$, of $H_T^*(W)$ over $\C[\lambda]$.
In this basis, $\mu$ acts by the first Hodge degree minus
$\dim_\C(W)/2$ and does not differentiate the coefficient $\lambda$.

\subsection{The calibrated grading}

Define
\begin{equation}
 \begin{split}
 G_d&=z\partial_z+\lambda\partial_\lambda+\tfrac12+\mu+\frac{c\cup}{z},\\
 G_a&=z\partial_z+\lambda\partial_\lambda+\tfrac12+\mu+\frac{c\star}{z}.
 \end{split}
 \label{shift:gradings}
\end{equation}
The subscripts distinguish the constant, or descendant, frame from the
quantum, or ancestor, frame.  Both operators act on coefficients in
$\lambda$, whereas the Novikov variables are held fixed.

\begin{lemma}
\label{shift:grading-lemma}
With $S_W=S_W(0,z)$ in the convention fixed above,
\begin{equation}
                 G_a=S_WG_dS_W^{-1}.
 \label{shift:grading-calibration}
\end{equation}
\end{lemma}

\begin{proof}
Let $D_c$ be the derivation of the Novikov ring defined by
\[
 D_c(Q^\beta y^l)
   =\left((r+1)l+\int_\beta\kappa\right)Q^\beta y^l.
\]
Homogeneity of the two-point invariants defining $S_W$ says
\begin{equation}
 (z\partial_z+\lambda\partial_\lambda+D_c)S_W
                      =S_W\mu-\mu S_W.
 \label{shift:homogeneity}
\end{equation}
This identity uses the first Hodge degree.  The evaluation maps are
algebraic morphisms, and the virtual class and cotangent-line classes
are algebraic.  Their push-pull operations therefore have the Hodge
bidegrees prescribed by virtual dimension, even when the inserted
classes are not algebraic.  Concretely, if the input and output basis
vectors have first Hodge degrees $h_j$ and $h_i$, the coefficient of
$Q^\beta y^l z^{-k-1}$ in the corresponding entry of $S_W$ has
$\lambda$-degree
\[
 h_j+k+1-h_i-\left((r+1)l+\int_\beta\kappa\right).
\]
Since $\lambda$ and $z$ both have bidegree $(1,1)$ for this calculation,
this is precisely \eqref{shift:homogeneity}.  It applies equally to
primitive and odd inputs.

The divisor equation gives
\begin{equation}
             zD_cS_W=(c\star)S_W-S_W(c\cup).
 \label{shift:divisor}
\end{equation}
The scalar equivariant part of $c$ occurs identically in the two
multiplications, so cancels in their difference.  Substituting
\eqref{shift:divisor} into \eqref{shift:homogeneity} gives
$G_aS_W=S_WG_d$.  These identities first hold in the expansion at
$z=\infty$.  Coefficientwise rationality of $S_W$, established by the
localization calculation, also gives the identities in its Laurent
expansion at $z=0$.
\end{proof}

\subsection{The shift operator and its fixed-locus formula}

Here is the genus-zero input from Iritani that we use.  For a smooth
projective variety $Y$ with a $\C^*$-action, form its associated bundle
$\widehat Y\to\mathbb P^1$.  Let $\sigma_{\min}$ be the section class
of the fixed component with positive normal weights.  The operator
defined by two-fiber section invariants, with monomials
$Q^{\widehat d-\sigma_{\min}}$, is an unlocalized equivariant
cohomology operator.  After reversing the sign of Iritani's loop
variable, it translates $\lambda$ to $\lambda+z$ and satisfies
\begin{equation}
 T_a=S_YT_dS_Y^{-1},\qquad
 T_d\big|_{F}
 =Q^{\sigma_F-\sigma_{\min}}
   \prod_{m,x}
   \frac{\prod_{a=-\infty}^{0}(x+m\lambda-az)}
        {\prod_{a=-\infty}^{-m}(x+m\lambda-az)}
    e^{z\partial_\lambda}.
 \label{shift:iritani-input}
\end{equation}
Here $x$ runs through the ordinary Chern roots of the weight-$m$
normal subbundle of $F$; each quotient has only finitely many remaining
factors.  This is the projective specialization of
\cite[Proposition~2.2, Remark~3.4, Definition~3.9,
Remark~3.10, Definition~3.13, and
Theorem~3.14]{IritaniShift}.  The source's seminegativity condition is
automatic for a complete target, since its global functions are
constant.  Its calibration $M$ is related to ours by
$S_Y(z)=M(-z)^{-1}$.

We describe the associated geometry in this application to fix both
the section classes and the signs.  For the action that has weight one
on the trivial summand of $E\oplus\mathcal O$, it is
\begin{equation}
 \widehat W
  =\mathbb P_{B\times\mathbb P^1}
     \bigl(\pr_B^*E\oplus\pr_{\mathbb P^1}^*\mathcal O(-1)\bigr).
 \label{shift:associated-space}
\end{equation}
Equivalently it is
$W\times(\C^2\setminus\{0\})/\C^*$, with
$s\cdot(w,v_1,v_2)=(s\cdot w,s^{-1}v_1,s^{-1}v_2)$.
An additional torus scales $v_2$, and its equivariant parameter is
Iritani's loop variable before sign reversal.  The fibers over
$[1:0]$ and $[0:1]$ carry the original action and the action composed
with this additional torus.  Identifying their equivariant
cohomologies by the induced change of torus variables accounts for
the translation in \eqref{shift:iritani-input}.

More explicitly, write $\zeta$ for that parameter before sign
reversal, and let $\rho_0(t,u)w=t\cdot w$ and
$\rho_1(t,u)w=(tu)\cdot w$ be the actions on the two fibers.  The
change of torus variables induces
$\Phi_1:H^*_{T\times\C^*,\rho_0}(W)\to
H^*_{T\times\C^*,\rho_1}(W)$ with
\[
 \Phi_1(f(\lambda,\zeta)a)
       =f(\lambda+\zeta,\zeta)\Phi_1(a).
\]
If $\iota_0,\iota_\infty$ denote the fiber inclusions, the section
correspondence is defined by
\begin{equation}
 (\widetilde T a,b)_\infty
  =\sum_{\beta,l}Q^\beta y^l
     \left\langle\iota_{0*}a,\iota_{\infty*}b
       \right\rangle^{\widehat W,T\times\C^*}
       _{0,2,\sigma_{\min}+(\beta,l)}.
 \label{shift:section-definition}
\end{equation}
Only effective section classes are included.  The inputs belong to
the equivariant cohomologies of their respective fibers, and the
pairing on the left is the one on the second fiber.  The operator
used here is $T_a=(\Phi_1^{-1}\widetilde T)|_{\zeta=-z}$.
Its semilinearity translates $\lambda$ to $\lambda+z$, as asserted.

In this construction the normal weight at $X$ is positive, so $X$
gives $\sigma_{\min}$.  A section obtained from a point of $X$ has
tautological degree zero in \eqref{shift:associated-space}; a section
obtained from the trivial-summand fixed component $B$ is the line
$\mathcal O(-1)$ and has tautological degree one.  Both have base
class zero.  Thus
$\sigma_B-\sigma_{\min}$ is precisely the fiber-line class of $W$.

This also checks the curve labels in the shift theorem.  The integral
projective-bundle splitting identifies $H_2(W,\Z)$ by
$(\beta,l)$.  The same splitting on \eqref{shift:associated-space}
identifies a section class by $(\beta,1,l)$.  Subtraction of
$\sigma_{\min}=(0,1,0)$ leaves $(\beta,0,l)$.  The inclusions of
the two fibers preserve $\beta$ and tautological degree, hence induce
the same identification on integral homology.  At every splitting
of a section stable map the labels therefore add in actual integral
homology.  In particular the use of \eqref{shift:iritani-input}
does not replace classes by their numerical equivalence classes.

\begin{proposition}
\label{shift:operator}
There is a shift operator
\[
 T_a=A(z,\lambda,Q,y)e^{z\partial_\lambda}
\]
on the small quantum-cohomology module of $W$, whose matrix $A$ is
polynomial in $z,\lambda$ at every curve coefficient in the chosen
equivariant basis.  Its calibration is
\begin{equation}
                   T_a=S_WT_dS_W^{-1},
 \label{shift:calibration}
\end{equation}
where, writing $v=x+w_F\lambda$ for the equivariant normal roots,
\begin{equation}
 T_d\big|_F
   =y^{j_F}
      \prod_x
       \begin{cases}
        x+\lambda,& F=X,\\[2pt]
        (x-\lambda-z)^{-1},& F=B
       \end{cases}
       e^{z\partial_\lambda},
 \qquad j_X=0,\quad j_B=1.
 \label{shift:fixed-formula}
\end{equation}
All operators act on the full equivariant cohomology with its super
pairing.
\end{proposition}

\begin{proof}
The target $W$ is smooth and projective, hence satisfies the
semiprojectivity and weight hypotheses of the cited theorem.
Its equivariant formality is also explicit in the basis
$\pi^*b\,p^j$.  Formula \eqref{shift:associated-space} is projective,
so its section invariants are defined by proper equivariant
pushforward.  The two fiber insertions are polynomial equivariant
classes, and the equivariant pairing of $W$ is perfect over
$\C[\lambda]$.  Consequently the operator is polynomial in both
equivariant parameters at each curve coefficient; changing the sign
of the second parameter and identifying the fibers preserves this
property.

The normal weights are $1$ at $X$ and $-1$ at $B$.  In
\eqref{shift:iritani-input} these give the factors $v$ and
$(v-z)^{-1}$, respectively.  The section calculation above gives
$y^{j_F}$.  Finally, in the convention for the fundamental solution
used here, its pairing formula and symplectic identity give
$S_W(z)=M(z)^*=M(-z)^{-1}$.  Iritani's identity
$M T_a=T_d M$, with his $z$ replaced by $-z$, is therefore exactly
\eqref{shift:calibration}.

The two-marked section correspondence and its localization proof
use arbitrary evaluation classes and diagonal contractions.  They
thus apply on full cohomology with the categorical inverse pairing
and its Koszul signs.  At primary background zero no assertion
about a power-series extension in odd background coordinates is
needed.
\end{proof}

\begin{lemma}
\label{shift:polynomiality}
For every fixed actual base class $\beta$, the matrix coefficients
of $A$ and of $c\star$ are polynomials in $y,z,\lambda$ (with no
$z$ dependence in $c\star$).  In particular they are holomorphic
functions of $y$ before any localization in $\lambda$.
\end{lemma}

\begin{proof}
The dual of the vector bundle in \eqref{shift:associated-space} is
the sum of the globally generated bundles $\pr_B^*E^*$ and
$\pr_{\mathbb P^1}^*\mathcal O(1)$.  Its tautological line bundle
is therefore globally generated.  Every effective section class has
$l\geq0$, and subtracting $\sigma_{\min}$ does not change $l$.
The same lower bound for $W$ was established by the master-space
construction.

Fix basis vectors for an input and a paired output.  The section
invariant defining the corresponding shift entry has two fiber
incidence insertions of fixed cohomological degrees.  The
projective-bundle formula on $\widehat W$ gives
\begin{equation}
 \operatorname{vdim}_{\C}
 \overline{\mathcal M}_{0,2}
       \bigl(\widehat W,\sigma_{\min}+(\beta,l)\bigr)
   =\dim_\C W+1+(r+1)l+\int_\beta\kappa.
 \label{shift:section-dimension}
\end{equation}
Thus virtual dimension increases by $r+1$ when $l$ increases by
one.  Proper equivariant integration is polynomial in the
parameters.  If virtual dimension exceeds the total degree of the
insertions, its result would have negative degree and is zero.
Thus only finitely many $l$ occur for this entry and $\beta$.
The finite polynomial inverse-pairing matrix used to recover
operator entries does not change this conclusion.  The identical
argument with the three primary insertions defining a quantum
product proves it for $c\star$.
\end{proof}

We have now obtained operators whose coefficients can be studied
both as Novikov series at $y=0$ and as functions of a nonzero complex
variable $y$.  Before making that change of viewpoint, we verify
that the shift respects the grading and determine its limiting
spectrum.

\begin{lemma}
\label{shift:commutation-lemma}
The grading and shift commute:
\begin{equation}
                          [G_a,T_a]=0.
 \label{shift:commutation}
\end{equation}
\end{lemma}

\begin{proof}
By the two calibration identities it suffices to work in the
descendant frame.  Equivariant homogeneity of restriction identifies
$\lambda\partial_\lambda+\mu$ on the summand for $F$ with
$\lambda\partial_\lambda+\mu_F-n_F/2$.  Also
\[
 c|_F=c_1(TF)+\sum_x(x+w_F\lambda).
\]
Write $T_d|_F=a_F E_z$, where $E_z=e^{z\partial_\lambda}$ and
$a_F$ is the multiplier in \eqref{shift:fixed-formula}.  The operator
$z\partial_z+\lambda\partial_\lambda$ commutes with $E_z$.
Combining it with first Hodge degree counts $1$ for each factor
$x+\lambda$ and $-1$ for each factor $(x-\lambda-z)^{-1}$.
The resulting commutator with $a_F$ is $W_F a_F$, where
$W_F=n_Fw_F$.  The power $y^{j_F}$ is held fixed in this calculation.
On the other hand,
\[
 E_z(c|_F)=(c|_F+W_Fz)E_z,
 \qquad
 [c|_F/z,a_FE_z]=-W_Fa_FE_z.
\]
The two contributions cancel.  Scalar degree-centering terms and
the ordinary cup multiplications in this computation introduce no
further commutators.
\end{proof}

\subsection{The spectrum of the fiber theory}

Let $Q^{>0}$ denote the ideal of positive base degree in the Novikov
ring.  Since an effective nonzero base class has positive ample
degree, reduction modulo this ideal retains exactly the maps with
constant projection to $B$.  The variable $y$ is retained.  We next
identify $T_a$ modulo $z$ in this fiber theory.

\begin{lemma}
\label{shift:fiber-operator}
On $H_T^*(W)$ with the fiber quantum product,
\begin{equation}
               T_a\bmod(z,Q^{>0})=(p+\lambda)\star.
 \label{shift:seidel}
\end{equation}
\end{lemma}

\begin{proof}
The product is linear over the classical superalgebra $H^*(B)$:
all evaluations of a genus-zero fiber map have the same projection
to $B$, so base classes pull out of the corresponding pushforward
with the usual signs.  The same reasoning applies to the two-point
operator $S_W$ in this reduction.  Formula
\eqref{shift:calibration}, whose shift fixes base cohomology, then
shows that $T_a$ is $H^*(B)$-linear.

Put $D_y=y\partial_y$.  The divisor equation gives
\[
 S_W(zD_y-p\cup)S_W^{-1}=zD_y-p\star.
\]
The operator on the left before conjugation commutes with $T_d$.
On $X$, the multiplier has no $y$ factor and $p|_X$ is independent
of $\lambda$.  On $B$, $p|_B=-\lambda$ and
\[
 [zD_y,T_d|_B]=zT_d|_B,
 \qquad [\lambda,T_d|_B]=-zT_d|_B.
\]
Thus the two terms cancel there as well.  Conjugating and reducing
modulo $z$ proves that $T_a|_{z=0}$ commutes with $p\star$.

Assign complex cohomological degree one to $p,\lambda,z$ and
degree $r+1$ to $y$ for the present fiber calculation.  Polynomiality
in $\lambda$ implies that the quantum powers of $p$ through $p^r$
are the classical powers: a positive $y$ contribution has degree
at least $r+1$.  These powers form a basis over $H^*(B)[\lambda]$,
so an $H^*(B)[\lambda]$-linear operator commuting with $p\star$ is
determined by its value on $1$.

In this same grading $T_a$ has degree one.  To see this directly
from the section definition, the normal bundle to a minimal
section has one summand $\mathcal O(-1)$ and the other summands
are trivial.  Hence $c_1(T\widehat W)\cdot\sigma_{\min}=1$.
Equivalently, \eqref{shift:section-dimension} at $\beta=0$ has
virtual dimension $\dim W+1+(r+1)l$.  Each fiber inclusion raises
insertion degree by one; comparing the resulting degree with the
fiber pairing gives operator degree $1-(r+1)l$ in the coefficient
of $y^l$.  It follows that $T_a(1)|_{z=0}$ has no
positive $y$ coefficient.  At $y=0$ all curve classes are zero,
$S_W=1$, and \eqref{shift:fixed-formula} applied to $1$ has
restrictions $p+\lambda$ on $X$ and zero on $B$.  These are exactly
the restrictions of $p+\lambda$ on $W$.  Thus
$T_a(1)|_{z=0}=p+\lambda$, which proves \eqref{shift:seidel}.
\end{proof}

Let $I=H^{>0}(B;\C)$.  This is a nilpotent ideal, also when
$H^*(B)$ has odd classes.  Reducing the fiber quantum algebra
modulo $I$ gives
\begin{equation}
 \C[\lambda,y,p]\big/\bigl(p^r(p+\lambda)-y\bigr).
 \label{shift:fiber-relation}
\end{equation}
Indeed the classical projective-bundle relation reduces to
$p^r(p+\lambda)=0$.  In degree $r+1$ the only possible positive
curve correction is a constant multiple of $y$.  Its coefficient
is the degree-one line invariant in a fiber $\mathbb P^r$, hence
one.  More explicitly, one may pair the relevant structure
constant with a top-degree base class; the degree-zero base
projection reduces the computation to the fiber.  The base
directions contribute no obstruction because
$H^1(C,\mathcal O_C)=0$ for a genus-zero domain.  The coefficient
has degree zero, so is independent of $\lambda$ and may equally
be computed at $\lambda=0$.  It is then
\[
       \left\langle H,H^r,H^r\right\rangle^{\mathbb P^r}
                      _{0,3,\mathrm{line}}=1.
\]
Here $H$ is the ordinary hyperplane class of the fiber.
Indeed two general point conditions determine a unique line, and
the remaining marked point is its unique intersection with a
general hyperplane.  Convexity of projective space makes this the
virtual count as well.  This gives the coefficient of $y$ in
$p\star p^r$, hence the coefficient in
\eqref{shift:fiber-relation}.

\begin{proposition}
\label{shift:spectrum-proposition}
For generic $(y,\lambda)$, the distinct eigenvalue branches of
$T_a\bmod(z,Q^{>0})$ are
\begin{equation}
             \lambda+h,\qquad h^r(h+\lambda)=y.
 \label{shift:spectrum}
\end{equation}
Each branch may have multiplicity and a nontrivial nilpotent part.
For $\lambda\ne0$, near $y=0$ there is one branch tending to zero
and a cluster of $r$ branches tending to $\lambda$.
\end{proposition}

\begin{proof}
On the quotient by $I$, \eqref{shift:seidel} and
\eqref{shift:fiber-relation} give the stated eigenvalues.  The
finite filtration by powers of $I$ is preserved by the operator,
and the induced operator on each successive quotient is the same
fiber multiplication tensored with $I^m/I^{m+1}$.  Its characteristic
polynomial therefore has the same set of roots.  Away from the
critical values of $h\mapsto h^r(h+\lambda)$ these roots are
distinct.  At $y=0$ the roots of that polynomial are $-\lambda$
and $0$, with multiplicities one and $r$, respectively; translation
by $\lambda$ gives the final assertion.  This argument uses
nilpotence of positive-degree base cohomology and does not impose
semisimplicity on its quantum cohomology.
\end{proof}

The polynomial dependence of the two operators and the finite set
of eigenvalue branches are the ingredients needed for the
spectral construction in the next section.  The calibration
identities retain their full curve labels, so the later return to
$y=0$ will still compare the individual descendant theories of
the two fixed manifolds.
